\section{Marco Teórico}
\label{sec:marco_teorico}

\subsection{Memoria Asociativa Entrópica}
\label{sec:eam}

La Memoria Asociativa Entrópica (EAM) \cite{pineda2021entropic} es un modelo de memoria
que opera sobre funciones discretas.
Una \textbf{pista} es una función discreta $f : \{1, \ldots, n\} \to \{0, \ldots, m-1\}$,
representada como un vector $\mathbf{a} \in \{0, \ldots, m-1\}^n$.
La dimensión $n$ es el número de \emph{características} y $m$ es el número de valores posibles
por característica (el alfabeto).

En la formulación original de 2021 el registro es \emph{booleano}: la memoria almacena
una función característica $R : A \times V \to \{0, 1\}$ y el registro se realiza por
disyunción lógica.
La variante que este trabajo utiliza es la \textbf{EAM ponderada} (W-EAM), introducida
en el aprendizaje fonético \cite{pineda2022weighted} y empleada desde entonces en toda
la línea de trabajo del grupo \cite{morales2022manuscript,pineda2023imagery}: la memoria
se representa como una relación ponderada $\mathbf{R} \in \mathbb{N}^{n \times m}$,
donde $\mathbf{R}[i, j]$ cuenta cuántas veces la característica $i$ ha tomado el valor
$j$ a lo largo de todos los registros.
Esta distinción importa para el lector que verifique contra las fuentes: las
definiciones que siguen corresponden a la variante ponderada, no a la booleana de 2021.

\subsubsection*{Operación de registro ($\lambda$)}

Dada una pista $\mathbf{a}$, se construye el patrón unitario:
\begin{equation}
    r(\mathbf{a})[i, j] = \begin{cases} 1 & \text{si } j = a_i \\ 0 & \text{en otro caso} \end{cases}
    \label{eq:patron-unitario}
\end{equation}
y se acumula en la memoria:
\begin{equation}
    \mathbf{R} \leftarrow \mathbf{R} + r(\mathbf{a})
    \label{eq:registro}
\end{equation}
La acumulación está acotada superiormente (suma con tope $w$) para evitar desbordamiento,
tal como definen las fuentes \cite{pineda2022weighted,morales2025missing}; los autores
describen esta operación como una forma de la regla de aprendizaje de Hebb.

La idea central es esta: la memoria \emph{no almacena
patrones individuales} sino que acumula una \textbf{distribución de experiencias}.
Cada columna $\mathbf{R}[i, \cdot]$ se convierte progresivamente en una distribución de
frecuencias sobre los valores posibles de la característica $i$.
La memoria va formando abstracciones y prototipos a partir de esta superposición
estadística; no se imponen desde el exterior ni se calculan de forma separada.

\subsubsection*{Operación de reconocimiento ($\eta$)}

Se dice que la pista $\mathbf{a}$ está \emph{contenida} en $\mathbf{R}$ si, para toda
característica $i$, la celda correspondiente tiene masa positiva: $\mathbf{R}[i, a_i] > 0$.
La versión con tolerancia $\xi$ admite a lo más $\xi$ fallas:
\begin{equation}
    \eta(\mathbf{a}, \mathbf{R}, \xi) = \mathbf{1}\!\left[\,\bigl|\{i : \mathbf{R}[i, a_i] = 0\}\bigr| \leq \xi\,\right]
    \label{eq:reconocimiento}
\end{equation}
Es una \emph{implicación material relajada}: la pista se acepta si cada una de sus
celdas marcadas encuentra soporte en la memoria, tolerando hasta $\xi$ columnas sin él.

Durante el reconocimiento se computan además los \textbf{pesos por dimensión}:
\begin{equation}
    w_i = \mathbf{R}[i,\, a_i]
    \label{eq:pesos}
\end{equation}
Estos pesos reflejan cuán frecuentemente la característica $i$ ha tenido el valor $a_i$
en los patrones registrados.
Una precisión de notación: en las definiciones formales de las fuentes, $\eta$ devuelve
un \emph{booleano} (aceptar/rechazar), y los $w_i$ son cantidades intermedias calculadas
durante la prueba de aceptación.
La implementación de referencia de Pineda \& Morales sí los expone como valor de
retorno de primera clase (\texttt{recog\_weights}), y este trabajo los usa de esa
manera: son la interfaz entre la memoria homo-asociativa y la hetero-asociativa
(véase la Sección~\ref{sec:hetero4d}).
Es una diferencia de empaquetado, no de sustancia.

El parámetro $\kappa$ define un umbral de reconocimiento relativo: se considera
reconocido si el peso medio $\bar{w} \geq \kappa \cdot \bar{\mu}$, donde $\bar{\mu}$ es
la media global de la memoria.
El parámetro $\iota$ controla la poda: en la \emph{relación operativa} de la memoria,
cualquier celda con peso $< \iota \cdot (\text{suma}/\text{celdas no cero})$ se
apaga temporalmente durante la recuperación, eliminando ruido estadístico de baja
frecuencia sin modificar la relación acumulada $\mathbf{R}$.

\subsubsection*{Operación de recuperación ($\beta$)}

La recuperación de un patrón a partir de una pista $\mathbf{a}$ procede columna a columna.
Para cada característica $i$, se construye una distribución de probabilidad sobre $\{0,\ldots,m-1\}$
modulada por una gaussiana de desviación $\sigma \cdot m$ centrada en el valor $a_i$
(cuando se dispone de la pista):
\begin{equation}
    p_i(j) \propto \mathbf{R}[i, j] \cdot \exp\!\left(-\frac{(j - a_i)^2}{2(\sigma m)^2}\right)
    \label{eq:distribucion-recuperacion}
\end{equation}
El valor recuperado $\hat{a}_i$ se obtiene \textbf{muestreando} esta distribución:
$\hat{a}_i \sim p_i(\cdot)$.
Esta modulación gaussiana corresponde al mecanismo de recuperación vigente desde el
trabajo de imaginería \cite{pineda2023imagery} (el artículo base de 2021 empleaba
distribuciones triangulares, hoy superadas); es la versión implementada en el código de
referencia que este sistema utiliza.

Esta variación no es un error del modelo; es parte de cómo recuerda:
\emph{recordar es muestrear}, no leer un valor determinista.
Diferentes ejecuciones de $\beta$ sobre la misma pista y la misma memoria pueden producir
resultados distintos.
El parámetro $\sigma$ controla la dispersión del muestreo: $\sigma = 0$ concentra la
masa en la pista exacta, valores mayores permiten explorar vecindades del espacio de
representación; valores altos favorecen recuperaciones imaginativas al costo de
falsas recolecciones \cite{pineda2023imagery}.

La \textbf{entropía} de la memoria se calcula como la media de las entropías de Shannon
por columna:
\begin{equation}
    H(\mathbf{R}) = \frac{1}{n} \sum_{i=1}^{n} \left(-\sum_{j=0}^{m-1} \hat{p}_{ij} \log_2 \hat{p}_{ij}\right)
    \label{eq:entropia}
\end{equation}
donde $\hat{p}_{ij} = \mathbf{R}[i,j] / \sum_l \mathbf{R}[i,l]$.
La entropía resume el estado de la memoria: una memoria con entropía alta ha acumulado
más variedad; una con entropía baja está dominada por pocos valores.
Conviene precisar que la entropía baja \emph{no} es intrínsecamente patológica: en las
fuentes, una función totalmente determinada tiene entropía exactamente cero y es una
memoria perfectamente bien formada, solo que rígida y reproductiva
\cite{pineda2021entropic,morales2022manuscript}.
Si la entropía baja es deseable o no depende de la tarea: un directorio que debe señalar
sin ambigüedad a un especialista \emph{quiere} baja ambigüedad, mientras que una memoria
de contenido que debe generalizar quiere un rango medio.
La evidencia empírica de las fuentes ubica el óptimo de reconocimiento y recuperación
constructiva en un rango \emph{intermedio} de entropía: baja entropía da precisión alta
con recall bajo, alta entropía lo contrario \cite{pineda2023imagery}.

\subsection{Memoria Hetero-asociativa 4D}
\label{sec:hetero4d}

La Memoria Hetero-Asociativa Entrópica (EHAM) de Morales y
Pineda~\cite{morales2024hetero,morales2025missing}
extiende la EAM al almacenamiento de pares de objetos posiblemente pertenecientes a dominios
o modalidades distintas.
En lugar de una relación bidimensional, la EHAM almacena las asociaciones en una relación
ponderada de cuatro dimensiones, de modo que una pista en un dominio puede proyectar una
relación en el dominio asociado.
Formalmente, relaciona un dominio izquierdo con $n$ características y alfabeto de tamaño $m$,
y un dominio derecho con $p$ características y alfabeto de tamaño $q$.
La relación se almacena en un tensor
$\mathbf{R} \in \mathbb{N}^{n \times p \times (m+1) \times (q+1)}$,
donde los tamaños $(m+1)$ y $(q+1)$ incorporan \emph{centinelas} en la posición cero
para representar funciones parciales (características no definidas).
Las operaciones $\lambda$, $\eta$ y $H$ se generalizan de forma directa: el registro es
la misma suma acotada sobre el tensor 4D, el reconocimiento es la misma implicación
material relajada con umbrales $\iota$ y $\kappa$, y la entropía es el promedio de las
entropías de Shannon por par de características $(a_i, b_j)$ \cite{morales2025missing}.

La recuperación hetero-asociativa introduce además el \textbf{problema de la pista faltante}
(\emph{missing cue problem}), contribución central del artículo publicado
\cite{morales2025missing}: cuando la pista fuente selecciona un plano de memoria en el
dominio destino pero no existe información adicional de ese dominio, el plano queda
ampliamente indeterminado y no hay con qué completar la recuperación.
Los autores proponen los métodos \emph{random-and-test}, \emph{sample-and-test} y
\emph{search-and-test} para ese caso límite.
Conviene ubicarlo con precisión: la proyección ponderada que se describe a continuación
es el mecanismo \emph{general} de recuperación cuando sí existe una pista fuente con
pesos; las soluciones a la pista faltante son el parche para el caso específico en que
ni siquiera esa información existe.
El presente sistema usa la EHAM como sustrato de asociación texto$\leftrightarrow$imagen
y como base para los directorios hetero-asociativos agente$\leftrightarrow$especialidad.

\subsubsection*{Registro hetero-asociativo}

Dada una pista izquierda $\mathbf{a} \in \{0,\ldots,m-1\}^n$ y una pista derecha
$\mathbf{b} \in \{0,\ldots,q-1\}^p$, el patrón unitario conjunto es:
\begin{equation}
    r(\mathbf{a}, \mathbf{b})[i, j, a_i, b_j] = 1
    \label{eq:registro-hetero}
\end{equation}
y se acumula de forma idéntica a la ecuación~\eqref{eq:registro}.

\subsubsection*{Proyección AND-ponderada}

Sea $\mathbf{a}$ una pista izquierda con pesos por dimensión $\mathbf{w}$ provenientes de
la operación $\eta$ de la memoria homo-asociativa del dominio izquierdo
(véase la ecuación~\eqref{eq:pesos}).
Se normalizan los pesos:
\begin{equation}
    w'_i = \frac{n \cdot w_i}{\displaystyle\sum_{k=1}^{n} w_k}
    \label{eq:normalizacion-pesos}
\end{equation}
de modo que $\frac{1}{n}\sum_i w'_i = 1$ (media unitaria).
La \textbf{proyección} $\Pi(\mathbf{a})$ sobre el dominio derecho se define por:
\begin{equation}
    \Pi[j, l] = \begin{cases}
        0 & \text{si } \exists\, i : \mathbf{R}[i, j, a_i, l] = 0 \\[4pt]
        \displaystyle\sum_{i=1}^{n} w'_i\, \mathbf{R}[i, j, a_i, l] & \text{en otro caso}
    \end{cases}
    \label{eq:proyeccion-AND}
\end{equation}
La condición de anulación es la \textbf{conjunción AND}: una celda del resultado se
anula si \emph{cualquiera} de las $n$ evidencias aporta cero.
Esto implementa \emph{containment} estricto: la proyección solo produce masa donde
todas las características de la pista tienen soporte simultáneo en la memoria.
Los pesos $w'_i$ ponderan la suma por la frecuencia con que cada característica de la pista
ha sido observada en la memoria homo-asociativa: características muy frecuentes en la pista
pesan más, lo cual permite que la memoria homo-asociativa module la hetero-asociativa
de forma nativa y justifica la arquitectura de cinco memorias por agente.

Sobre la atribución de esta regla: el nombre ``proyección AND-ponderada'' es de este
trabajo, pero el mecanismo no es una síntesis propia; es exactamente la operación
\texttt{project()} de la implementación de referencia de Pineda \& Morales (que este
sistema incorpora sin modificaciones), donde la anulación conjuntiva y la suma ponderada
por pesos normalizados son líneas literales de esa única función, operando sobre la
relación ya podada por $\iota$.
Tanto la recuperación $\beta$ del artículo \cite{morales2025missing} como las lecturas
de directorio de este sistema se construyen sobre esa primitiva compartida.

La recuperación hetero-asociativa $\beta$ aplica muestreo por columna sobre $\Pi$ y,
cuando es necesario, una búsqueda local (\emph{sample-and-search}) para mejorar la
calidad del patrón recuperado.

\subsubsection*{Asimetría entre dominios}

La memoria homo-asociativa de un dominio modela la \emph{distribución marginal} de ese
dominio: genera pesos $w_i$ que reflejan cuán típica es una pista dentro de ese espacio.
La memoria hetero-asociativa, en cambio, modela la \emph{distribución conjunta} de dos
dominios y, por construcción, \textbf{no genera pesos por característica} del dominio derecho.
Esta asimetría es propia del modelo y justifica que cada agente en EAM-TMS requiera
\emph{cinco} memorias: dos homo-asociativas (una por dominio, como
árbitros de pertenencia y fuente de pesos), una hetero-asociativa de contenido, y
dos hetero-asociativas de directorio (una de texto y una visual, una por modalidad de pista).

\subsection{Memoria Transactiva}
\label{sec:tmt}

El modelo de memoria transactiva de Wegner \cite{wegner1987transactive} describe cómo
los grupos desarrollan sistemas de memoria distribuida cuya capacidad excede la suma
de sus partes.
El mecanismo central es la separación entre \emph{contenido} y \emph{directorio}:
mientras cada miembro almacena información especializada, todos los miembros mantienen
un \emph{metadirectorio} de quién sabe qué.
Cabe señalar que este puente entre la teoría de Wegner y las memorias asociativas
entrópicas es una contribución de este proyecto: los artículos de la línea EAM/EHAM no
contienen nociones de directorio, agentes ni memoria transactiva.

Los tres procesos fundamentales son:

\begin{enumerate}
    \item \textbf{Actualización del directorio} (\emph{directory updating}):
    tras cada interacción de recuperación exitosa, \emph{todos} los miembros del grupo,
    incluso quienes no respondieron, actualizan su directorio anotando quién
    fue el especialista.
    Este proceso distribuye la información de especialización incluso a los miembros que
    no participaron activamente en la respuesta.

    \item \textbf{Asignación de información} (\emph{information allocation}):
    la información nueva se encamina automáticamente hacia el miembro cuya especialidad
    mejor corresponde.
    En la fase temprana este enrutamiento puede requerir un mediador que evalúe las
    capacidades de cada miembro; en la fase madura el propio directorio del emisor
    basta para determinarlo.

    \item \textbf{Coordinación de recuperación} (\emph{retrieval coordination}):
    cuando un miembro recibe una consulta fuera de su dominio, consulta su directorio
    y reenvía la consulta al especialista correspondiente sin intervención de ningún
    coordinador central.
    El \textbf{rechazo claro}, cuando ningún miembro reconoce el dominio de la
    consulta, es un comportamiento válido del grupo y evita que el sistema invente
    especialistas inexistentes.
\end{enumerate}

En la \textbf{fase temprana}, el grupo depende de un mediador: los miembros
no conocen aún las especialidades de los demás y las consultas deben ser difundidas a
todos.
La \textbf{fase madura} se alcanza cuando cada miembro ha acumulado suficiente evidencia
en su directorio para enrutar directamente: el mediador sale del circuito y las
interacciones se vuelven punto a punto.

La \textbf{entropía del directorio} resume el estado del sistema transactivo:
un directorio con entropía máxima $\log_2 K$ (siendo $K$ el número de agentes;
$\log_2 8 = 3$ bits en esta versión) indica especialización perfectamente balanceada;
una entropía baja indica colapso hacia un solo agente o formación incompleta.
A diferencia de la entropía de una memoria de contenido, aquí el diagnóstico se hace
sobre la \emph{distribución de registros entre agentes}, y el estado deseable es el
máximo: que el grupo reparta la experticia de forma pareja.
